\section{Stability above a strict barrier}
\label{sec:big-reference-stability}

Retain the arbitrary smooth big background $\theta$, the strict potential
$\rho\leq0$, and $\Omega=X\setminus\{\rho=-\infty\}$ from
\cref{sec:big-minimal-comparison}. The same auxiliary-subsolution argument
as in \cref{sec:bounded-stability} applies to a reference which is itself
unbounded. The reference-preserving construction and the domination
principle above replace their bounded counterparts. This reconstructs
\cite[Theorem~3.3 and Remark~3.4]{Dang26} in the minimal case, with the
normalization of the comparison potential explicit.

\begin{theorem}\label{thm:big-reference-one-sided-stability}
Fix $C_0,L>0$, $p>1$, and
\[
 0<\gamma<\frac{p-1}{np+p-1}.
\]
Let $\eta\in\PSH(X,\theta)$ satisfy $\rho\leq\eta\leq0$. Suppose
$u\in\PSH(X,\theta)$ satisfies
\[
 \eta-C_0\leq u\leq\eta,\qquad
 \mu_u=f\,\mathrm{d}V,\qquad 0\leq f\in L^p(\mathrm{d}V),
 \quad\|f\|_p\leq L.
\]
There is a constant $C$, depending only on the fixed geometry,
$\rho,C_0,L,p,\gamma$, such that every $v\in\PSH(X,\theta)$ with
$v\leq\eta$ satisfies
\begin{equation}\label{eq:big-reference-stability}
 \sup_\Omega(v-u)_+
 \leq C\left(\int_\Omega(v-u)_+\,\mathrm{d}V\right)^\gamma.
\end{equation}
The constant is independent of $\eta$ in the stated range. The comparison
potential $v$ need not have the ordinary singularity type of $\eta$.
\end{theorem}

Differences in this statement are taken on $\Omega$, where $u$ and $\eta$
are locally bounded; their integrals are unchanged by extending them by
zero on the analytic complement. The bound also means the corresponding
order inequality for the canonical potentials on all of $X$.

\begin{proof}
Replace $v$ by $\max(v,u)$, so that
$\eta-C_0\leq u\leq v\leq\eta$ and $0\leq v-u\leq C_0$ on $\Omega$.
Fix $q\in(1,p)$ and $\delta>0$, and put
\[
 E_\delta=\{u<v-\delta\}\cap\Omega,\qquad
 g_\delta=\mathbf1_{E_\delta}f,\qquad m_\delta=\|g_\delta\|_q.
\]
When $m_\delta>0$, \cref{lem:big-reference-subsolution}, followed by
subtracting one from its potential, gives $h\in\PSH(X,\theta)$ and a
fixed $b>0$ such that
\[
 \eta-1\leq h\leq\eta,\qquad
 \mu_h\geq b\,g_\delta/m_\delta\,\mathrm{d}V.
\]
The constant is uniform in $\eta$ because the auxiliary density has
$L^q$ norm one and the strict barrier is fixed.

If $0<m_\delta<b/2$, set
\[
 t=(2m_\delta/b)^{1/n},\qquad
 w=(1-t)v+th-\delta-C_0t.
\]
Since $h-v\leq C_0$, the set $\{u<w\}\cap\Omega$ is contained in
$E_\delta$. On this set the local positive-product expansion gives
$\mu_w\geq t^n\mu_h\geq2\mu_u$. Both potentials are bounded below
by $\rho$ minus a constant, and $w\leq u+C_0$. Thus
\cref{thm:big-barrier-domination} applies and gives $u\geq w$. It follows
that
\[
 v-u\leq\delta+C_0t+t(v-h)\leq\delta+(C_0+1)t.
\]
If $m_\delta=0$, apply the same domination theorem with coefficient zero
to $u$ and $v-\delta$. If $m_\delta\geq b/2$, use $v-u\leq C_0$.
In all cases,
\[
 \sup_\Omega(v-u)\leq\delta+C_1m_\delta^{1/n}.
\]
Write $I=\int_\Omega(v-u)\,\mathrm{d}V$. For $I=0$, the canonical
potentials agree; otherwise H\"older and Chebyshev give
\[
 m_\delta\leq L(I/\delta)^{(p-q)/(pq)}.
\]
For $0<I\leq1$, choosing
$\delta=I^{(p-q)/(npq+p-q)}$ yields the exponent
$\gamma_q=(p-q)/(npq+p-q)$. Choose $q>1$ close enough to one that
$\gamma_q>\gamma$. The case $I>1$ follows from the bound by $C_0$.
This proves \eqref{eq:big-reference-stability}.
\end{proof}

\begin{corollary}\label{cor:big-reference-uniform-convergence}
Fix $\rho\leq\eta\leq0$, $C_0,L>0$, and $p>1$. Consider the potentials
$u\in\PSH(X,\theta)$ satisfying
\[
 \eta-C_0\leq u\leq\eta,\qquad
 \mu_u=f_u\,\mathrm{d}V,\qquad0\leq f_u\in L^p(\mathrm{d}V),
 \quad\|f_u\|_p\leq L.
\]
This family is compact for uniform convergence of differences on $\Omega$.
If $u_j\to u$ in $L^1$, then $u$ belongs to the same family and
\[
 \sup_\Omega|u_j-u|\leq C\|u_j-u\|_{L^1(\mathrm{d}V)}^\gamma
 \longrightarrow0
\]
for every exponent in \cref{thm:big-reference-one-sided-stability}.
Moreover $f_{u_j}$ converges weakly to $f_u$ in $L^p$.
\end{corollary}
\begin{proof}
Apply the theorem in both orders to $u_j,u_k$, with the common reference
$\eta$. It follows that an $L^1$-Cauchy sequence is uniformly Cauchy in
differences on $\Omega$. The uniform limit there is locally
$\theta$-plurisubharmonic and agrees almost everywhere, hence everywhere
in canonical representatives, with the $L^1$ limit $u$. The bounds
$\eta-C_0\leq u\leq\eta$ hold on all of $X$. Letting $k\to\infty$ in
the pairwise estimate proves the displayed quantitative bound, before
using any equation for $u$.

The densities form a bounded set in the reflexive space $L^p$. For a
weakly convergent subsequence, say $f_{u_{j_k}}\rightharpoonup f$, local
uniform convergence of the bounded potentials gives
$\mu_u=f\,\mathrm{d}V$ on $\Omega$, by Bedford--Taylor continuity for
the smooth nonclosed background. Both measures give zero mass to the
analytic complement, so this equality holds on $X$. Positivity and
weak lower semicontinuity give $f\geq0$ and $\|f\|_p\leq L$.
Every weak cluster point is therefore the same density $f_u$, proving
weak convergence of the whole sequence. Finally, the common bounds by
$\eta$ and $\eta-C_0$ give $L^1$ compactness by quasi-plurisubharmonic
compactness. The preceding argument upgrades each convergent subsequence
to uniform convergence in the stated family.
\end{proof}

For a normalized solution with minimal singularities, take
$\eta=V_\theta$. Every normalized comparison potential lies below this
reference. The compactness assertion uses a common reference and a
uniform $L^p$ density bound; it makes no claim from weak convergence of
arbitrary non-pluripolar measures. No metric mass property is imposed.

\section{Local H\"older continuity in a big background}
\label{sec:big-minimal-holder}

We now reconstruct the local H\"older conclusion of
\cite[Theorem~A]{Dang26}. Two points require attention when passing from
\cref{sec:holder-from-stability}: the solution can be unbounded on $X$,
and the smooth lower bound for its Hessian need not be closed. We retain
Demailly's classical local regularization expansion and explain the
quantitative estimate used here.

\begin{lemma}\label{lem:smooth-background-regularization-rate}
Let $a$ be quasi-plurisubharmonic on $X$, with $\sup_Xa=0$ and
$\theta+\ddc a\geq0$, and suppose that $a$ is locally bounded on $\Omega$.
There are constants $A,K,C>0$ such that the averaging operators $R_t$
from \cref{lem:chern-kiselman-input} satisfy
\[
 R_ta+Kt^2\downarrow a,\qquad
 \|R_ta-a\|_{L^1(\mathrm{d}V)}\leq Ct^2.
\]
For small $\delta>0$ and $c>0$, the transform on $\Omega$
\[
 \Phi_{c,\delta}
 =\inf_{0<t\leq\delta}
 \{R_ta+K(t^2-\delta^2)+K(t-\delta)-c\log(t/\delta)\}
\]
extends canonically to $X$ and satisfies
\begin{equation}\label{eq:smooth-background-transform-rate}
 a-K(\delta^2+\delta)\leq\Phi_{c,\delta}\leq R_\delta a\leq0,
 \qquad \theta+\ddc\Phi_{c,\delta}\geq-(Ac+C\delta)\omega.
\end{equation}
At each point of $\Omega$ the infimum is attained at a positive radius.
On every compact $K_1\Subset\Omega$, for sufficiently small $t$,
$\|\nabla R_ta\|_{K_1}\leq C_{K_1}/t$.
\end{lemma}

\begin{proof}
Apply \cite[Theorem~4.1]{Dem94} to the closed current $\ddc a$ and the
smooth lower bound $-\theta$. It gives monotonicity and logarithmic
convexity, with
$\lambda(x,t)=t\partial_t(R_ta+K_0t^2)\geq0$ increasing in $t$.
We first justify a linear error for this smooth lower bound, rather than
infer a rate from the unspecified modulus in that theorem.

Choose the auxiliary curvature form to be $A\omega$, with $A$ large
enough that the Chern curvature plus $A\omega\otimes\operatorname{Id}$
is nonnegative on every decomposable tensor. Equivalently, after
absorbing the normalization constants, its coefficients satisfy
\[
 \sum R_{j\bar k\ell\bar m}\xi_j\bar\xi_k\zeta_\ell\bar\zeta_m
 \geq-A|\xi|^2|\zeta|^2.
\]
This choice is available by compactness. Pairing this inequality with
the nonnegative eigenvalues of a positive complex Hessian bounds its
curvature contraction by $-A$ times its trace. Hence the algebraic inequality used in
\cite[Lemma~4.4]{Dem94} holds with no small $\epsilon$ term and no
$M_\epsilon$ correction.

Fix $x$ and $t=|w|$. In a uniformly sized coordinate chart about $x$,
write $\beta=\ddc|z|^2$ and
choose a quadratic $q_{x,t}$ with
\[
 \ddc q_{x,t}=\theta(x)+Bt\beta,
\]
where $B$ is fixed and large enough that $a+q_{x,t}$ is plurisubharmonic
on a ball containing the averaging support and its nearby base points.
This follows from the smooth variation bound
$\|\theta(y)-\theta(x)\|\leq C|y-x|$ on that ball of radius $O(t)$.
For the Hessian calculation at $(x,w)$ the quadratic is held fixed.
Its coefficients and all needed derivatives are uniformly bounded.

Here are the error estimates in the classical expansion
\cite[Proposition~3.8, equations (3.11), (4.3), and (4.5)]{Dem94}.
Normalized quasi-plurisubharmonic compactness bounds $\|a\|_1$.
On each fixed larger chart, add one fixed strictly plurisubharmonic
quadratic $Q$ to $a$. If
$\nu=\ddc(a+Q)\wedge\beta^{n-1}$, the local Lelong mass inequality and
a fixed larger-ball mass bound give
$\nu(B_t(x))\leq Ct^{2n-2}$. The coefficients of
$\partial\bar\partial(a+q_{x,t})$ are controlled in total variation
by those of $\ddc(a+Q)$ plus a uniformly bounded smooth form. Scaling
therefore gives, uniformly in $x$ and $t$,

\[
 \int_{|\zeta|<1}|\partial\bar\partial(a+q_{x,t})|
       (\operatorname{exph}_x(t\zeta))\,\mathrm{d}\zeta
 \leq Ct^{-2}.
\]
The integral is understood as the pullback of coefficient measures for
nonsmooth $a$. Thus the terms with a coefficient $O(t^3)$ contribute
$O(t)$; the order-two expansion has an $O(t)$ remainder as well.
The mixed and vertical contributions have bounds
$C(|\xi||\tau|+|\tau|^2)$. The curvature term is bounded below by
$-A\lambda(x,t)|\xi|^2$, using the preceding curvature choice.
Equation (4.5) identifies its radial mass with $\lambda$ up to $O(t^2)$.
Subtracting the regularization of $q_{x,t}$ restores $-\theta(x)$;
its base Hessian differs from $\ddc q_{x,t}$ by $O(t)$, its mixed and
vertical Hessians are bounded, and its radial logarithmic derivative
is $O(t^2)$. The extra $Bt\beta$ also costs only $O(t)$.
Consequently the full base--fiber Hessian satisfies
\[
 \ddc_{x,w}R_{|w|}a[\xi,\tau]^2
 \geq-\theta_x[\xi]^2-(A\lambda(x,t)+Ct)|\xi|^2
       -K_1(|\xi||\tau|+|\tau|^2).
\]
All constants are independent of $x,t,c,\delta$ and uniform in the
normalized family. In particular, setting $\xi=0$ in the same local
expansion gives a uniform lower bound for the fiber Hessian. Adding a
fixed multiple of $|w|^2$ gives its subharmonicity, hence the
logarithmic convexity and monotonicity above with a uniform $K_0$.

Increase $K\geq K_0$ and add $K(t^2+t)$ to the regularization. The
positive vertical Hessian of $K|w|$ absorbs the mixed term by
\[
 K_1|\xi||\tau|\leq C't|\xi|^2+\frac{K}{8t}|\tau|^2,
\]
and absorbs the remaining bounded vertical error for small $t$.
The logarithmic derivative of the corrected function is at least
$\lambda$. The capped Kiselman minimum argument in
\cref{lem:general-kiselman-modulus} therefore gives the last inequality
in \eqref{eq:smooth-background-transform-rate}, now with error $C\delta$.
On compact subsets of $\Omega$, local boundedness of $a$ makes the
logarithmic penalty diverge uniformly at zero, so that the minimum
occurs away from zero. The same argument handles minimizers at the
upper cap without differentiating their radius.

The other bounds in \eqref{eq:smooth-background-transform-rate} follow
from monotonicity and evaluation at $t=\delta$. In particular the
transform is bounded above near the analytic complement. The removable
singularity theorem extends it, preserving the exact smooth lower
Hessian bound by local freezing and passage to the limit. Its lower
bound by $a-K(\delta^2+\delta)$ also extends by canonical representatives.

The adjoint computation in \cref{lem:chern-smoothing-integral} gives
$R_t^*1=1+O(t^2)$. It needs only $a\in L^1$, since
\[
 \left|\int_X(R_ta-a)\,\mathrm{d}V\right|
 \leq Ct^2\|a\|_1.
\]
Together with $R_ta-a\geq-Kt^2$, this proves the stated $L^1$ estimate.
Finally the smoothing kernel near $K_1$ samples a fixed larger compact
subset of $\Omega$, where $a$ is bounded. The same kernel derivative
estimate as in \cref{lem:chern-smoothing-integral} gives $C_{K_1}/t$.
\end{proof}

\begin{theorem}\label{thm:big-minimal-holder}
Let $u\in\PSH(X,\theta)$ satisfy
\[
 \sup_Xu=0,\qquad V_\theta-C_0\leq u\leq V_\theta,
 \qquad\mu_u=f\,\mathrm{d}V,
 \quad0\leq f\in L^p(\mathrm{d}V),\quad p>1.
\]
Then $u$ is locally H\"older continuous on $\Omega$, with every exponent
\[
 0<\beta<\frac{2(p-1)}{np+3(p-1)}.
\]
For a fixed bound on $C_0$ and $\|f\|_p$, the constants are uniform on
each fixed compact subset of $\Omega$.
\end{theorem}

\begin{proof}
Subtract a constant from the strict barrier so that $\rho\leq u-1$.
It still satisfies $\theta+\ddc\rho\geq\kappa\omega$ for a fixed
$\kappa>0$. Choose
$0<\gamma<(p-1)/(np+p-1)$ and write $\alpha=2\gamma<1$.
For small $\delta$ set
\[
 c=\delta^\alpha,\qquad e=Ac+C\delta,\qquad s=e/\kappa,
 \qquad v_\delta=\frac{\Phi_{c,\delta}+s\rho}{1+s}.
\]
The preceding lemma gives $v_\delta\in\PSH(X,\theta)$. Its lower bound
for $\Phi_{c,\delta}$ implies $\rho\leq\Phi_{c,\delta}$ when $\delta$
is small. Thus
\[
 v_\delta\leq\Phi_{c,\delta}\leq R_\delta u\leq0.
\]
Apply \cref{thm:big-reference-one-sided-stability} with reference
$V_\theta$ and comparison $v_\delta$. The $L^1$ smoothing estimate gives
\[
 \sup_\Omega(v_\delta-u)_+
 \leq C\|(R_\delta u-u)_+\|_1^\gamma\leq C\delta^\alpha.
\]
Fix a compact subset on which $\rho\geq-M$. Since $s\leq C\delta^\alpha$
and $u\leq0$, the identity
\[
 \Phi_{c,\delta}-u=(1+s)(v_\delta-u)+s(u-\rho)
\]
gives $\Phi_{c,\delta}-u\leq C_M\delta^\alpha$ there.

At a minimizing radius $t$ the lower monotone bound yields
\[
 -c\log(t/\delta)\leq\Phi_{c,\delta}-u+K(\delta^2+\delta)
 \leq C_Mc.
\]
Hence $t\geq a_M\delta$ for a fixed $a_M>0$. As in the proof of
\cref{thm:hermitian-holder-elementary}, monotonicity now gives
\[
 R_{a_M\delta}u-u\leq C_M\delta^\alpha,
 \qquad R_{a_M\delta}u-u\geq-K(a_M\delta)^2.
\]
The constants are uniform on the chosen compact subset. Apply this to a
slightly larger compact neighborhood, where also
$\|\nabla R_tu\|\leq C_M/t$. For nearby points at distance $d$, smooth
approximation and a short path in that neighborhood give
\[
 |u(x)-u(y)|\leq C_M(t^\alpha+d/t).
\]
Take $t=d^{1/(1+\alpha)}$. This gives exponent
$\alpha/(1+\alpha)$; letting $\gamma$ approach its upper endpoint gives
the asserted range. Boundedness on the compact neighborhood handles
pairs farther apart.
\end{proof}

The density condition in this theorem is global. Smoothness of $f$ only
on $\Omega$ does not provide that hypothesis. The theorem gives no
continuity assertion across the pole set and does not assert the same
regularity for an arbitrary prescribed ordinary singularity type.
